% !TeX program = xelatex
\documentclass[11pt, a4paper]{article}

\usepackage[a4paper, margin=2.5cm]{geometry}
\usepackage{fontspec}
\usepackage[lithuanian, provide=*]{babel}
\babelprovide[import, onchar=ids fonts]{lithuanian}
\babelfont{rm}{Verdana}

\usepackage{amsmath}
\usepackage{amssymb}
\usepackage{booktabs}
\usepackage{enumitem}
\usepackage{tasks}

\setlist[itemize]{label=-}

\title{\textbf{Pasiruošimas perrašymui}}
\author{Uždavinių rinkinys}
\date{}

\begin{document}

\maketitle

\section*{I dalis. Užduotys}

\textbf{1.} Nustatykite duotų reiškinių tipus (R - racionalusis, T - trupmeninis, S - sveikas, I - iracionalusis).
\begin{tasks}(2)
    \task \(25 ^{-2}+ x ^{3} -x+ \sqrt{2}\)
    \task $ \frac{5}{x}+4x-7 $
    \task \(\sqrt{x-4}+2x+1\)
    \task \(\frac{x}{x^2-1}+ \sqrt{\pi}\)
    \task \(7x^{-1}+x+1\)
    \task \(\frac{x+8}{ \sqrt{2} } +x\)
\end{tasks}

\vspace{0.5cm}
\textbf{2.} Suprastinkite reiškinius ir nurodykite, ar gauti reiškiniai yra vienanariai. Jei reiškinys yra vienanaris, nurodykite jo koeficientą ir laipsnį; jei ne, paaiškinkite, kodėl.
\begin{tasks}(2)
    \task \(\displaystyle \frac{(-2x^5y^{-4})^{10} \cdot (8x^{-2}y^3)^{-15}}{(16x^4y^{-6})^{10}}\)
    \task \(\displaystyle \frac{(3x^2y^{-1})^4 \cdot (9x^{-2}y^2)^{-5}}{(3xy^{-3})^2}\)
    \task \(\displaystyle \frac{(5x^{-2}y^3)^8 \cdot (25x^3y^{-1})^{-10}}{(5x^{-7}y^2)^8}\)
    \task \(\displaystyle \frac{(4x^3y^{-2})^5 \cdot (2x^{-1}y^4)^{-9}}{8x^2y^{-5}}\)
    \task \(\displaystyle \frac{5x^6y^{-2}}{12a^{-3}b^{-4}} : \frac{25x^{-2}y^{-7}}{18a^4b^{-4}}\)
    \task \(\displaystyle \frac{14x^{-2}y^6}{21a^4b^{-3}} : \frac{2x^5y^{-2}}{9a^{-1}b^2}\)
    \task \(\displaystyle \frac{7x^3y^{-5}}{15a^{-2}b^6} : \frac{14x^{-4}y^2}{45a^3b^{-6}}\)
    \task \(\displaystyle \frac{9x^{-4}y^8}{10a^5b^{-2}} : \frac{3x^{-7}y^3}{25a^6b^{-2}}\)
\end{tasks}

\vspace{0.5cm}
\textbf{3.} Apskaičiuokite reiškinio reikšmę be smegenų tuštintuvo ir daugybos stulpeliu:
\begin{tasks}(2)
    \task \(\displaystyle \frac{72^3 + 28^3}{0,1^{-2}(72 \cdot 44 + 28^2)}\)
    \task \(\displaystyle \frac{\sqrt[3]{179^9}+121 \cdot\left(93^2+93 \cdot 56+28^2\right)}{(5 \sqrt{12}+\sqrt{75}-5 \sqrt{3})^2 \cdot\left(\sqrt{121^4}+179 \cdot 58\right)}\)
    \task \(\displaystyle \frac{(\sqrt{243}-\sqrt{111})(\sqrt{243}+\sqrt{111})\sqrt{132^4}-64 \cdot 13^3}{(4\sqrt{5})^2\cdot\left((2\sqrt{33})^4+11\cdot13\cdot48+52^2\right)}\)
\end{tasks}

\vspace{0.5cm}
\textbf{5.} Raskite kiekvieno reiškinio didžiausią ir mažiausią reikšmes (jei jos egzistuoja) ir nurodykite, su kokia $x$ reikšmei kiekviena reikšmė įgyjama.
\begin{tasks}(4)
    \task \(4x^2 + x - 1\)
    \task \(-x^2 + 3x + 1\)
    \task \(3x^2 + 2x + 2\)
    \task \(-2x^2 + 5x - 2\)
\end{tasks}

\vspace{0.5cm}
\newpage
\textbf{6.} Duotą reiškinį pertvarkykite ir užrašykite pavidalu \(a(x+b)^2+c\). Didelius reikškinius tvarkykite ir prastinkite principu skaldyk ir valdyk.
\begin{tasks}(1)
    \task \((2x+1)^3-x(6x+1)(x-2)-2\left(x^3+9x^2+3x\right)\)
    \task \((x-1)^2(1-x)-(x+1)^3-\left(-2x^3-3x^2-4x+1\right)\)
    \task \((1-2x)^3-8(x-1)(x+1)(1-x)-\left(3x^2-8x+9\right)\)
    \task \((x+2)^3-(3-x)^3+2(x-2)^2(x+1)-4\left(x^3-2x^2+8x-3\right)\)
\end{tasks}

\vspace{0.5cm}
\textbf{7.} Raskite parametrus \(a\), \(b\), \(c\), kad lygybė būtų teisinga su visomis \(x\) reikšmėmis:
\begin{tasks}(1)
    \task \((ax+b)(2x^2 - x + 3) + cx = 4x^3 - 6x^2 + 10x - 6\)
    \task \((ax-2b)(3x^2 + 2x - 1) + cx = 6x^3 + 10x^2 + 5x - 2\)
    \task \((2ax+b)(x^2 - 3x + 2) + cx = -2x^3 + 9x^2 - 8x + 6\)
    \task \((ax+3b)(-2x^2 + x + 4) + cx = -2x^3 + 13x^2 + 2x - 24\)
\end{tasks}

\vspace{0.5cm}
\textbf{8.} Raskite reiškinių mažiausias reikšmes.
\begin{tasks}(1)
    \task \(x^2 + 4x + y^2 - 8y + z^2 + 2z + 25\)
    \task \(x^2 - 6x + y^2 + 4y + z^2 - 2z + 20\)
    \task \(2x^2 + 8x + y^2 - 10y + 3z^2 + 6z + 38\)
    \task \(4x^2 - 8x + y^2 + 6y + z^2 + 4z + 25\)
\end{tasks}

\newpage

\section*{II dalis. Atsakymai}

\textbf{1.}
\begin{tasks}(2)
    \task R, S
    \task R, T
    \task I
    \task R, T
    \task R, T
    \task R, S
\end{tasks}

\vspace{0.5cm}
\textbf{2.}
\begin{tasks}(2)
    \task \(2^{-75}x^{40}y^{-25}\) – ne vienanaris, nes kintamojo \(y\) laipsnio rodiklis yra neigiamas.
    \task \(3^{-8}x^{16}y^{-8}\) – ne vienanaris, nes kintamojo \(y\) laipsnio rodiklis yra neigiamas.
    \task \(5^{-20}x^{10}y^{18}\) – vienanaris; koeficientas \(5^{-20}\), laipsnis \(28\).
    \task \(2^{-2}x^{22}y^{-41}\) – ne vienanaris, nes kintamojo \(y\) laipsnio rodiklis yra neigiamas.
    \task \(\frac{3}{10}a^7x^8y^5\) – vienanaris; koeficientas \(\frac{3}{10}\), laipsnis \(20\).
    \task \(3a^{-5}b^5x^{-7}y^8\) – ne vienanaris, nes kintamųjų \(a\) ir \(x\) laipsnių rodikliai yra neigiami.
    \task \(\frac{3}{2}a^5b^{-12}x^7y^{-7}\) – ne vienanaris, nes kintamųjų \(b\) ir \(y\) laipsnių rodikliai yra neigiami.
    \task \(\frac{15}{2}ax^3y^5\) – vienanaris; koeficientas \(\frac{15}{2}\), laipsnis \(9\).
\end{tasks}

\vspace{0.5cm}
\textbf{3.}
\begin{tasks}(3)
    \task $1$
    \task $1$
    \task $1$
\end{tasks}

\vspace{0.5cm}
\textbf{5.}
\begin{tasks}(2)
    \task Mažiausia reikšmė: \(-\frac{17}{16}\), kai \(x=-\frac{1}{8}\); didžiausios reikšmės nėra.
    \task Didžiausia reikšmė: \(\frac{13}{4}\), kai \(x=\frac{3}{2}\); mažiausios reikšmės nėra.
    \task Mažiausia reikšmė: \(\frac{5}{3}\), kai \(x=-\frac{1}{3}\); didžiausios reikšmės nėra.
    \task Didžiausia reikšmė: \(\frac{9}{8}\), kai \(x=\frac{5}{4}\); mažiausios reikšmės nėra.
\end{tasks}

\vspace{0.5cm}
\textbf{6.}
\begin{tasks}(2)
    \task \(5(x+0,2)^2+0,8\)
    \task \(3\left(x-\frac{1}{3}\right)^2-\frac{4}{3}\)
    \task \((x-3)^2-9\)
    \task \(-(x-3,5)^2+13,25\)
\end{tasks}

\vspace{0.5cm}
\textbf{7.}
\begin{tasks}(1)
    \task \(a = 2\), \(b = -2\), \(c = 2\)
    \task \(a = 2\), \(b = -1\), \(c = 3\)
    \task \(a = -1\), \(b = 3\), \(c = 5\)
    \task \(a = 1\), \(b = -2\), \(c = 4\)
\end{tasks}

\vspace{0.5cm}
\textbf{8.}
\begin{tasks}(4)
    \task $4$
    \task $6$
    \task $2$
    \task $8$
\end{tasks}

\end{document}